\documentclass[../../../include/open-logic-section]{subfiles}

\begin{document}
\olfileid{sth}{card-arithmetic}{expotough}

\olsection[一些简化]{基数指数运算的一些简化}

虽然定义 $\canonord$ 略显复杂，但由此得到了一个关于基数加法与乘法的有用结论，即 \olref[simp]{cardplustimesmax}。然而，超限指数运算却无法如此直接地加以简化。为了说明原因，我们首先给出一条将有限情形下常见模式加以推广的命题（尽管其证明的抽象程度较高）：

\begin{prop}\ollabel{simplecardexpo}
对于任意基数 $\cardfont{a}, \cardfont{b}, \cardfont{c}$，有 $\cardexpo{\cardfont{a}}{\cardfont{b} \cardplus \cardfont{c}} =
\cardexpo{\cardfont{a}}{\cardfont{b}} \cardtimes
\cardexpo{\cardfont{a}}{\cardfont{c}}$ 以及
$\cardexpo{(\cardexpo{\cardfont{a}}{\cardfont{b}})}{\cardfont{c}} =
\cardexpo{\cardfont{a}}{\cardfont{b} \cardtimes \cardfont{c}}$。
\end{prop}

\begin{proof}
对于第一个论断，考虑一个函数 $f \colon
(\cardfont{b}\disjointsum\cardfont{c}) \to \cardfont{a}$。现在将其“拆分”：对每个 $\beta \in \cardfont{b}$ 定义 $f_{\cardfont{b}}(\beta) = f(\beta, 0)$，对每个 $\gamma \in \cardfont{c}$ 定义 $f_{\cardfont{c}}(\gamma) = f(\gamma, 1)$。映射 $f \mapsto
(f_{\cardfont{b}} \times f_{\cardfont{c}})$ 是一个双射 $\funfromto{\cardfont{b} \disjointsum \cardfont{c}}{\cardfont{a}} \to
(\funfromto{\cardfont{b}}{\cardfont{a}} \times
\funfromto{\cardfont{c}}{\cardfont{a}})$。

对于第二个论断，考虑一个函数 $f \colon \cardfont{c} \to
(\funfromto{\cardfont{b}}{\cardfont{a}})$；因此对每个 $\gamma \in \cardfont{c}$，便得到某个函数 $f(\gamma) \colon \cardfont{b} \to
\cardfont{a}$。现在对每个 $\tuple{\beta, \gamma} \in \cardfont{b} \times \cardfont{c}$ 定义 $f^*(\beta, \gamma) = (f(\gamma))(\beta)$。
映射 $f \mapsto f^*$ 是一个双射 $\funfromto{\cardfont{c}}{(\funfromto{\cardfont{b}}{\cardfont{a}})}
\to \funfromto{\cardfont{b} \cardtimes \cardfont{c}}{\cardfont{a}}$。
\end{proof}

现在，我们\emph{希望}的是，在处理无穷基数时，能有一种简便的方法来计算 $\cardexpo{\cardfont{a}}{\cardfont{b}}$。这是朝此方向迈出的很好的一步：

\begin{prop}\ollabel{cardexpo2reduct}
如果 $2 \leq \cardfont{a} \leq \cardfont{b}$ 且 $\cardfont{b}$ 是无穷的，那么 $\cardexpo{\cardfont{a}}{\cardfont{b}} =
\cardexpo{2}{\cardfont{b}}$
\end{prop}

\begin{proof}
\begin{align*}
	\cardexpo{2}{\cardfont{b}} &\leq
	\cardexpo{\cardfont{a}}{\cardfont{b}}\text{，因为 $2 \leq \cardfont{a}$}\\
	&\leq \cardexpo{(2^{\cardfont{a}})}{\cardfont{b}}
	\text{，由 \olref[opps]{lem:SizePowerset2Exp}}\\
	&= \cardexpo{2}{\cardfont{a} \cardtimes \cardfont{b}}
	\text{，由 \olref{simplecardexpo}} \\
	&= \cardexpo{2}{\cardfont{b}}
	\text{，由 \olref[simp]{cardplustimesmax}}
\end{align*}
\end{proof}

我们确实不应指望能对此作进一步简化，因为由 \olref[card-arithmetic][opps]{lem:SizePowerset2Exp} 可知 $\cardfont{b} < \cardexpo{2}{\cardfont{b}}$。然而，这并未告诉我们当 $\cardfont{b} < \cardfont{a}$ 时，对 $\cardexpo{\cardfont{a}}{\cardfont{b}}$ 该作何论断。当然，如果 $\cardfont{b}$ 是\emph{有限的}，我们便知道该如何处理。

\begin{prop}
如果 $\cardfont{a}$ 是无穷的且 $n \in \omega$，则 $\cardexpo{\cardfont{a}}{n} = \cardfont{a}$
\end{prop}

\begin{proof}
$\cardexpo{\cardfont{a}}{n} = \cardfont{a} \cardtimes \cardfont{a}
\cardtimes \ldots \cardtimes \cardfont{a} = \cardfont{a}$，由 \olref[simp]{cardplustimesmax}。
\end{proof}
\noindent
此外，在其他一些情形下，我们可以控制 $\cardexpo{\cardfont{a}}{\cardfont{b}}$ 的大小：

\begin{prop}
如果 $2 \leq \cardfont{b} < \cardfont{a} \leq
\cardexpo{2}{\cardfont{b}}$ 且 $\cardfont{b}$ 是无穷的，则 $\cardexpo{\cardfont{a}}{\cardfont{b}} = \cardexpo{2}{\cardfont{b}}$
\end{prop}

\begin{proof}
$\cardexpo{2}{\cardfont{b}}\leq \cardexpo{\cardfont{a}}{\cardfont{b}} \leq \cardexpo{(\cardexpo{2}{\cardfont{b}})}{\cardfont{b}} = \cardexpo{2}{\cardfont{b}\cardtimes\cardfont{b}} = \cardexpo{2}{\cardfont{b}}$，推理过程同 \olref{cardexpo2reduct}。
\end{proof}
\noindent
但是，超过这一点，情况就变得颇为微妙了。

\end{document}
